\section{Discussion}
\label{sec.discussion}

\paragraph{Confirmation rate}

The pruning depth grows with the round. The fraction of the chain it removes
does not. For a chain of height $\rate fr$, \Cref{cor.limit} gives
$\prn(r)/\hgt\to\advr(1+\advr-\rate)/\rate$, which does not depend on $r$. At
$\rate=1$ the fraction is $\advr^2$. So a client that receives an honest chain
at round $r$ keeps a prefix of height about $(1-\advr^2)fr$, and one that
receives an honest chain at a later round keeps strictly more. The confirmed
height grows at rate $(1-\advr^2)f$. Confirmation does not deteriorate as the
chain ages. Only the depth measured in blocks does, and
\Cref{cor.noconstant} shows that it must.

\paragraph{Idempotence}

\Cref{def.client} is applied to a chain as it is received, never to an output
of a previous application. Applying it twice at the same round is sound and
lossy. It is sound because the output is a chain of $\bt_r$ like any other, so
\Cref{cor.round} applies to it. It is lossy because, in the normalised form of
\Cref{cor.limit}, a second application sends $\rate$ to
$(1+\advr)(\rate-\advr)$. This map has the single fixed point $1+\advr$ and
slope $1+\advr>1$ there, so the fixed point repels and every
$\rate<1+\advr$ falls below $\advr$ after enough applications, at which point
the client outputs $\bot$. Nothing is gained by the repetition:
\Cref{cor.prune} confirms each output for every later round, so an output
never has to be pruned again.

\paragraph{Client knowledge}

\Cref{def.client} uses $n$, $t$, $f$, $\epsilon$ and $\lambda$, and the client
is given them. A client that prunes more than $\prn(r)$ still outputs a prefix
of every honest chain, since a prefix of a prefix is a prefix. A client that
prunes less has no guarantee from \Cref{cor.round}. An error in the parameters
that raises the pruning depth therefore costs liveness, and one that lowers it
costs safety. We do not analyse how a client obtains the parameters.

\paragraph{Variable difficulty}

The analysis assumes a static population and a fixed target $T$. Under these
assumptions the height of a chain is a proxy for the work in it, which is what
\Cref{eq.early} needs. Variable difficulty breaks the proxy. An adversary that
controls the timestamps of the blocks it produces can inflate the apparent time
between difficulty adjustments, induce the largest decrease in target the rule
allows, and reach a given height for much less work.
\cite{bitvm3} reports that about a fifth of the hashing power suffices to keep
pace with the honest validators in this way, against an on-chain client that
checks height alone. It repairs the client by binding the committed timestamps
to the consensus time of the chain. \Cref{def.client} reads height against the round,
so it is open to the same substitution. A version for variable difficulty would
have to compare accumulated work against the round rather than height against
the round. We do not carry that out here.

\paragraph{Open questions}

The additive terms of \Cref{thm.prefix,cor.round} are of order $\lambda f$,
the order of the common prefix parameter of~\cite{gkl24}. We do not claim
that they are optimal. The analysis is for the synchronous network
of~\cite{gkl24}. In the bounded delay model of~\cite{gkl24,pss16} the rate
$f$ is replaced by a rate that depends on the delay $\Delta$, and we expect
$t/(n-2t)$ to be replaced by the corresponding ratio of rates.
